\begin{frame}[fragile]{두 단계 전이: ``다른'' 작업에서 배운 특징}
  실전 전이학습은 \textbf{다른(보통 훨씬 큰) 작업}에서 배운
  가중치를 가져온다: (1) \textbf{큰} 데이터셋(4,000장,
  period 3 줄무늬 --- ``ImageNet''의 위장)에서
  \texttt{SmallCNN} 학습(100\%) $\to$
  (2) 합성곱 가중치를 \texttt{copy\_} 로 복사, 합성곱
  \textbf{동결}하고 \textbf{작은} 데이터셋(300장, period 4 ---
  ``새 문제'')의 \textbf{헤드만} 20 에폭 학습 $\to$
  \textbf{검증 정확도 100\%}.
\begin{lstlisting}
pre = SmallCNN()                       # 1단계: 큰 데이터에서 사전학습
# ... 15 에폭 학습 ...   # 학습 정확도 100%

model = SmallCNN()
with torch.no_grad():                  # 2단계: 가중치 복사
    for d, s in zip(model.conv1.parameters(),
                    pre.conv1.parameters()):
        d.copy_(s)
    # ... conv2도 동일하게 ...
for p in model.conv1.parameters():
    p.requires_grad = False            # 합성곱 동결
# ... conv2도 ...
opt = torch.optim.Adam(
    filter(lambda p: p.requires_grad,
           model.parameters()), lr=0.01)
# ... 20 에폭 ...   # 검증 정확도 100%
\end{lstlisting}
  \vskip0.2em
  \kq{중요한 디테일: 새 작업의 줄무늬 \textbf{주파수가 변했다}
  (period 3 $\to$ 4).} 사전학습된 특징은 새 작업의 ``정답
  특징''이 \textbf{아니다} --- 그런데도 헤드 재학습만으로
  100\%. 왜? 첫 합성곱 층이 배운 것은 ``줄무늬의 주기 3''이
  아니라 \textbf{``줄이 있다는 사실'', 즉 모서리·대비} 때문
  --- 주기 자체는 \textbf{뒤쪽 층(헤드)}이 담당. ``앞쪽 =
  범용, 뒤쪽 = 특화''라는 \textbf{비대칭이 직접 관찰}됨.
\end{frame}
